\begin{longtable}{p{29mm}p{72mm}p{55mm}}
\caption{Input and clinical-reference specifications, with unresolved source requirements.}\label{tab:input_reference_specification}\\
\toprule Element & Current analysis specification & Evidence boundary \\ \midrule\endfirsthead
\toprule Element & Current analysis specification & Evidence boundary \\ \midrule\endhead
EASI reference & Recorded total EASI, 0--72; observed outcomes only, with no imputed validation truth. & Regional-sum discrepancies are audited; the authoritative field in discrepant records and blinded repeat-assessment reliability are not confirmed. \\ \addlinespace
BSA input & Affected BSA entered from 0 to 100 (10\% is entered as 10). It supplies a clinical extent component of the hybrid estimate. & The source does not establish an observer-independent measurement protocol; this is not a blood-only score. \\ \addlinespace
Molecular inputs & CCL17, CCL26 and IL-5. Positive finite integration-export numerics recorded in pg/mL are log2 transformed once. & A positive numeric export is not proof of validated quantification. Transfer to another export or assay requires confirmed scale compatibility. \\ \addlinespace
Specimen and assay & Laboratory reports describe serum, Luminex LXSAHM-13/Luminex-200 and MSD K15067M-1/K15067M-2/SQ120 assays. & Matched specimen handling, technical-repeat precision and cross-platform bridging are not established by the analysis. \\ \addlinespace
Dilution and limits & Reports describe 50-microlitre serum and twofold laboratory dilution. No additional dilution factor is applied to the integration export. & Whether dilution is already corrected in that export, and validated lower/upper quantification limits, require source confirmation. Software calibration boundaries are not validated LLOQ/ULOQ. \\ \addlinespace
Clinical/serum timing & Patient identifier plus nominal W0, W12 or W24 links the clinical and molecular record. & Exact calendar-day separation and a permitted sampling window are unverified; nominal matching does not establish same-day measurement. \\ \addlinespace
Development scope & Age $\geq$12, observed EASI/BSA and all three positive molecular inputs. Fitting and selection use W0/W12 with patient-grouped inner and outer validation. & The marker panel arose from earlier exploration. W24 had previously been inspected and is same-cohort temporal application, not untouched external validation. \\ \addlinespace
Calculator output & The precise frozen coefficients return EASI units clipped to 0--72. Missing, nonfinite or nonpositive molecular inputs and invalid BSA are rejected. & BSA=0 returns zero by the area constraint; it does not prove absence of molecular inflammation. Independent validation of the single final coefficient vector is unavailable. \\ \addlinespace
Recorded IGA & Recorded grades 0--5 are retained. IGA 0/1 and decrease $\geq$2 use those recorded values without conversion to a 0--4 scale. & Exact instrument/version and correspondence to a validated IGA system require source confirmation. \\ \addlinespace
PP-NRS and DLQI & Recorded PP-NRS uses 0--10; adult DLQI uses 0--30 and age $\geq$18. A PP-NRS decrease $\geq$4 requires observed baseline $\geq$4 and observed follow-up. & PP-NRS recall window and exact questionnaire versions remain unconfirmed. Adult DLQI is unknown in adolescent composite rules. \\ \addlinespace
IgE and eosinophils & Analyses retain the recorded numeric values; adjusted models use log(1+x) and back-transformation. & Clinical laboratory units and the meaning of repeated IgE values of 5000 require source confirmation; no censoring threshold is inferred from repetition alone. \\ \addlinespace
Treatment targets & The September investigator framework supplies stage-specific physician and patient thresholds, with the stated three-valued missing-component logic. & The exact approved and prospectively implemented framework version and original voting records are not verified; target attainment is not validation of a new consensus or management effect. \\ \addlinespace
\bottomrule\end{longtable}
